Placing researchers in research groups that match their interests is difficult to perform manually and to scale, even though doing so well can support scientific collaboration, mentoring, and institutional planning~\cite{yang2014multilevel,sun2015raf}; as the number of researchers and groups grows, manual matching becomes increasingly subjective~\cite{zhang2023scholarly}. Scholarly recommender systems have addressed related tasks through content-based and interaction-based methods~\cite{zhang2023scholarly,bai2019scientific}, but interaction data such as co-authorship and citation links may be sparse for new researchers. Publication titles and abstracts instead provide direct evidence of current research interests and can support recommendations without prior interaction records~\cite{bai2019scientific,beel2016research}.

Topic modeling converts publication text into compact thematic representations using three common paradigms: probabilistic models such as LDA~\cite{jelodar2019lda,blei2003lda}, matrix-factorization models such as NMF~\cite{lee2000nmf}, and contextual embedding-based models such as BERTopic, which combines document embeddings, dimensionality reduction, density-based clustering, and class-based TF-IDF~\cite{grootendorst2022bertopic}. Previous studies have compared these families on social media, online communities, news, and other short-text collections~\cite{egger2022topic,kaur2024moving,babalola2024comprehensive}, but most have assessed topic interpretability or clustering behavior rather than the value of topic vectors inside a downstream recommendation task.

This creates two related gaps: existing studies have not established how LDA, NMF, and BERTopic compare under identical profile construction, similarity ranking, and researcher-level evaluation for researcher-to-research-group recommendation, and model selection has often relied on intrinsic topic coherence rather than this kind of controlled comparison, although prior work has questioned whether higher coherence predicts better downstream performance~\cite{roder2015exploring,rudiger2022topic}. These gaps motivate two research questions: (1) how do LDA, NMF, and BERTopic compare when used to rank research groups for a test researcher under one identical pipeline, and (2) does a topic model's intrinsic coherence agree with its downstream recommendation quality?

This study addresses these questions through a comparative evaluation of LDA, NMF, and BERTopic for publication-based researcher-to-research-group recommendation, evaluating 11,111 prepared publication documents from 3,257 researchers across 12 research groups using a researcher-level train-test split with group profiles built only from training researchers. As reported in Section~\ref{sec:results}, the answers diverge: the model with the most coherent topics is not the model with the best research-group rankings, indicating that coherence alone is not a reliable criterion for selecting a topic representation for this task. The main contributions are:
\begin{enumerate}
    \item A comparative evaluation of probabilistic, matrix-factorization, and contextual topic representations under the same profile construction, cosine-similarity ranking, and test protocol.
    \item A standardized researcher-level evaluation procedure in which test researchers remain separate from model fitting and group-profile construction.
    \item An analysis of the relationship between $C_v$ topic coherence and 13 recommendation metrics, including Precision, Recall, F1-score, NDCG, and MRR.
    \item A paired per-researcher significance analysis (Wilcoxon signed-rank tests) quantifying which observed ranking differences between models are statistically robust rather than attributable to sampling variation in the single split.
\end{enumerate}
